\chapter{The Saturation Regime: Multitone Harmonic Balance}
\label{chap:saturation}

\section{Why a new formulation is needed}

At large signal power the assumption underpinning \cref{chap:floquet} fails: the
generated sidebands carry enough power to modify the pump. Pump and signal must
then be found \emph{simultaneously} as a single nonlinear problem.

This is not a small change. The Floquet formulation of \cref{chap:floquet}
treats the sideband ladder as a linear system that can be factored once; that is
exactly what makes a gain map affordable. Solving pump and signal together
removes that structure --- the sidebands now appear inside the nonlinear
residual --- and restores the full Newton machinery of \cref{chap:pumphb} at a
much larger basis.

\begin{remark}[the cost is structural, not an implementation artefact]
The multitone Jacobian is block-dense in tone index, so memory scales as
$(n_{\mathrm{pump}} + 2S + 1)^{2}$ (\cref{chap:numerics}). On the same circuit,
moving from a gain-map pump basis with $H=10$ to a multitone basis at $S=10$
with $H=31$ takes the Jacobian from \num{2.46e6} to \num{23.6e6} nonzeros --- a
factor $9.6 \approx (31/10)^{2}$ --- while the dimension grows only $3.1\times$.
That factor is exactly the cost the Floquet split avoids, and it is the price of
asking a question the split cannot answer.
\end{remark}

\section{The tone lattice}

\subsection{Two incommensurate frequencies}

A pumped, signal-driven circuit has two fundamental frequencies. Rather than
treat the signal as a perturbation at $\ws$, we work with the pump frequency
$\wp$ and the \emph{detuning}
\begin{equation}
  \delta = \wp - \ws,
\end{equation}
and index every tone by an integer pair.

\begin{definition}[tone index]
A tone is a lattice coordinate $(h,q)\in\mathbb{Z}^{2}$ with physical frequency
\begin{equation}
  \boxed{\;\omega_{(h,q)} = h\,\wp + q\,\delta .\;}
  \label{eq:tone-frequency}
\end{equation}
\end{definition}

The three fundamental tones are
\begin{center}
\begin{tabular}{@{}l l l@{}}
\toprule
Tone & Coordinate & Frequency \\
\midrule
Pump   & $(1,0)$  & $\wp$ \\
Signal & $(1,-1)$ & $\wp - \delta = \ws$ \\
Idler  & $(1,1)$  & $\wp + \delta = 2\wp - \ws$ \\
\bottomrule
\end{tabular}
\end{center}
and the mixing relation $2\wp = \ws + \wi$ is the lattice identity
$2(1,0) = (1,-1) + (1,1)$ --- four-wave mixing becomes exact integer arithmetic.

\begin{remark}[why this indexing]
The alternative --- indexing by absolute frequency --- makes mixing a
floating-point coincidence test, which is fragile. On the lattice, whether two
tones mix is decided by integer addition, and the alias structure of the
transform is exact. The pump harmonics ($q=0$) and the signal sidebands
($q=\pm1$) live in the same index space, so a single basis carries both.
\end{remark}

\subsection{Positive-frequency representation}

Only positive-frequency tones are retained; a negative-frequency combination is
represented by the conjugate of its positive partner.

\begin{implnote}[\code{canonicalize} and the basis invariants]
\code{canonicalize} maps $(h,q)$ to a positive-frequency representative and a
conjugation flag. The \code{MultiToneBasis} constructor enforces:
\begin{itemize}
\item no DC tone $(0,0)$, and no exactly zero-frequency combination;
\item no duplicates;
\item no tone \emph{and} its conjugate both retained (that would double-count);
\item every retained tone strictly positive in frequency;
\item the three required tones $(1,0)$, $(1,-1)$, $(1,1)$ all present.
\end{itemize}
Each is checked with a specific error message. These are not defensive
boilerplate: a basis violating any of them produces a solvable but physically
meaningless system.
\end{implnote}

\subsection{The torus and its exact transform}

Two frequencies means a two-dimensional torus. Introduce phases
$\theta_p = \wp t$ and $\theta_\delta = \delta t$ and sample on an
$N_p\times N_\delta$ grid,
\begin{equation}
  \theta_p^{(r)} = \frac{2\pi r}{N_p},
  \qquad
  \theta_\delta^{(s)} = \frac{2\pi s}{N_\delta}.
\end{equation}
Synthesis is
\begin{equation}
  \boxed{\;
  \xvec(\theta_p,\theta_\delta)
  = 2\,\Real\!\left[\sum_{(h,q)\in\toneset}
      X_{(h,q)}\,e^{i(h\theta_p + q\theta_\delta)}\right],
  \;}
  \label{eq:torus-synthesis}
\end{equation}
the same positive-phasor convention as \cref{conv:phasor}, now in two
dimensions.

\begin{implnote}[FFT-based, exact]
\code{scatter\_signed\_spectrum} places each retained coefficient at
$(h \bmod N_p,\; q \bmod N_\delta)$ and its conjugate at the mirrored index;
\code{synthesize} is then a 2-D inverse FFT scaled by $N_pN_\delta$, and
\code{project} is the forward FFT divided by the same factor followed by a
gather of the retained modes. The transform pair is exact --- no interpolation,
no windowing.
\end{implnote}

\begin{proposition}[alias-free grid sizes]
Products of retained tones reach $(h+h', q+q')$. The grid must therefore satisfy
\begin{equation}
  N_p \ge 2\max_{h,h'}|h+h'| + 1,
  \qquad
  N_\delta \ge 2\max_{q,q'}|q+q'| + 1,
\end{equation}
with both rounded up to even.
\end{proposition}

\begin{implnote}[defaults are computed from the basis]
\code{MultiToneBasis.\_\_post\_init\_\_} computes these minima from the actual
tone list and uses them as defaults; explicitly supplied values are validated
against them and rejected if too small. A user cannot silently under-sample.
\end{implnote}

\subsection{Basis construction}

Two constructors are used in practice.

\paragraph{Sideband-matched (production).} Cover exactly the Floquet sidebands
$m\in[-S,S]$ of \cref{chap:floquet}, plus the pump harmonics. Each signed
sideband $m$ is first written as $(m+1,-1)$; negative-frequency results are
canonicalised, which naturally supplies the $q=+1$ representatives.

\begin{implnote}[coverage is proved, not assumed]
\code{build\_sideband\_matched\_basis} maps every requested sideband, then
verifies
\begin{lstlisting}
covered = basis.covered_sidebands()
if covered != set(range(-sidebands, sidebands + 1)):
    raise ValueError(f"sideband coverage mismatch: missing=..., extra=...")
\end{lstlisting}
with the inverse map
\begin{lstlisting}
def floquet_sideband_of(tone):
    if tone.q == -1: return tone.h - 1
    if tone.q ==  1: return -(tone.h + 1)
    return None
\end{lstlisting}
It also raises if \code{omega\_max} would clip a requested sideband or a pump
tone, and if two requests fold onto the same positive-frequency tone. A basis
that silently omits a sideband would produce a plausible, wrong compression
number.
\end{implnote}

\paragraph{Lattice.} A rectangular block $h\in\text{pump modes}$,
$q\in[-Q,Q]$, clipped by $\omega_{\max}$. Used by the convergence study, where
the natural knob is the signal order $Q$ rather than a Floquet sideband count.

\paragraph{Three-tone.} $\{(1,0),(1,-1),(1,1)\}$ only. Valid \emph{only} with a
fundamental-only pump basis.

\begin{caveat}[three-tone is a trap]
If the pump basis retains harmonics, a three-tone basis silently drops them ---
which changes the pump waveform and therefore every derived quantity. The driver
raises rather than dropping any pump $(h,0)$ mode.
\end{caveat}

\section{The multitone problem}

\subsection{Residual}

The structure is identical to \cref{chap:pumphb} with the mode index $k$
replaced by the tone index $(h,q)$:
\begin{equation}
  \boxed{\;
  R_{(h,q)}(\Xvec;\tau)
  = \Dmat\bigl(\omega_{(h,q)}\bigr)X_{(h,q)}
  + N_{(h,q)}(\Xvec)
  - S_{(h,q)}(\tau) = 0,
  \;}
\end{equation}
with $\Dmat(\omega)$ the same dynamic block \eqref{eq:dynamic-block} and the
nonlinear term computed by the same AFT cycle, now on the torus:
\begin{align}
  \psi &= \Bphi^{\mathsf T}\xvec(\theta_p,\theta_\delta) + \psi_{\mathrm{dc}},\\
  \iJ  &= \text{branch}(\psi) - \text{branch}(\psi_{\mathrm{dc}}),\\
  N &= \mathcal{P}\left[\Bphi\,\iJ\right],
\end{align}
where $\mathcal{P}$ is the torus projection. Note the same DC-current
subtraction as \cref{impl:dc-subtraction}, and the same delegation to the
circuit's branch law --- so a SNAIL device solves without any change to this
layer.

\subsection{The affine source path}

Continuation in \cref{chap:pumphb} scaled a single source from zero. Here there
are two sources, and the useful paths turn one on while holding the other fixed.

\begin{definition}[affine source path]
\begin{equation}
  \Svec(\tau) = \Svec_{\mathrm{start}} + \tau\,\Svec_{\Delta},
  \qquad \tau\in[0,1].
\end{equation}
\end{definition}

\begin{implnote}[three named paths]
\begin{center}
\begin{tabular}{@{}l l l@{}}
\toprule
Constructor & $\Svec_{\mathrm{start}}$ & $\Svec_\Delta$ \\
\midrule
\code{pump\_turn\_on} & $\vect{0}$ & $\Svec_p$ \\
\code{signal\_turn\_on} & $\vect{0}$ & $\Svec_p + \Svec_s$ \\
\code{signal\_substep} & $\Svec_p + \Svec_s^{\mathrm{old}}$
                       & $\Svec_s^{\mathrm{target}} - \Svec_s^{\mathrm{old}}$ \\
\bottomrule
\end{tabular}
\end{center}
The third is the important one: it holds the pump at full strength and moves
only the signal, which is the physically meaningful continuation when marching
up a compression curve. Because the tangent predictor uses
\code{source\_coeffs(1) - source\_coeffs(0)} (\cref{sec:continuation}), it gives
$\Svec_\Delta$ correctly for all three without modification.
\end{implnote}

\subsection{Preconditioning}

The preconditioner hierarchy transfers unchanged. The exact real-coupled
construction \eqref{eq:real-superblock} is assembled over tone pairs, with the
required $\Khat$ offsets being exactly $\{k-q\}\cup\{k+q\}$ over all retained
tone pairs --- now differences and sums of \emph{lattice coordinates}.

\begin{implnote}[cache survives \code{dataclasses.replace}]
A compression sweep rebuilds the problem once per signal-power point with only
\code{source\_path} changed. Everything expensive --- the dynamic blocks, the
scatter map, the symbolic factorisation --- depends on the circuit and basis
only. \code{FullMultiToneProblem} therefore carries an explicit \code{cache}
dict that survives \code{replace}, so that work happens once per sweep rather
than once per power point. Forgetting to carry \code{preconditioner} across the
rebuild silently reverts to the default; that is now explicit in
\code{\_problem\_with\_path}.
\end{implnote}

\section{Compression}

\subsection{Definition}

\begin{definition}[compression and $P_{1\mathrm{dB}}$]
For small-signal gain $G_0$ and gain $G(P_s)$ at input signal power $P_s$,
\begin{equation}
  \mathrm{compression}(P_s) = G_0 - G(P_s)\quad[\si{\dB}],
\end{equation}
and the input-referred 1\,dB compression point $P_{1\mathrm{dB}}$ is the
smallest $P_s$ at which the compression reaches \SI{1}{\dB}.
\end{definition}

\subsection{Four reference states}

Reporting gain and depletion in a normalised way requires four solves:

\begin{center}
\begin{tabular}{@{}l l l@{}}
\toprule
State & Pump & Signal \\
\midrule
\code{pump\_off\_signal\_on} & off & finite \\
\code{pump\_on\_signal\_off} & on & off \\
\code{pump\_on\_signal\_infinitesimal} & on & $\SI{1e-12}{\ampere}$ \\
\code{pump\_on\_signal\_finite} & on & finite \\
\bottomrule
\end{tabular}
\end{center}

From these,
\begin{align}
  \text{gain} &= 20\log_{10}\frac{|S_{21}^{\mathrm{signal}}|_{\text{pump on}}}
                                  {|S_{21}^{\mathrm{signal}}|_{\text{pump off}}},
  \label{eq:mt-gain}\\
  \text{pump depletion} &= 20\log_{10}
    \frac{|S_{21}^{\mathrm{pump}}|_{\text{signal finite}}}
         {|S_{21}^{\mathrm{pump}}|_{\text{signal off}}}.
  \label{eq:mt-depletion}
\end{align}
The infinitesimal-signal state exists to provide a well-conditioned seed for the
finite-signal solve and to give the small-signal gain in the same basis as the
compressed one --- so that the compression is a difference of like quantities.

\begin{caveat}[read gain through \code{tone\_s21}]
Gain must be read through \code{multitone.observables.tone\_s21} and reported as
the pump-on/pump-off ratio. An earlier path reconstructed it as
$|i\omega X|/V_{\mathrm{in}}$; that was biased by \SI{12.041}{\dB} on the SNAIL
benchmark. \code{tone\_s21} applies the reconstruction factor of two, forms
$V = 2i\omega X$, and routes through
\code{port\_s\_from\_unit\_current\_response}.
\end{caveat}

\subsection{\texorpdfstring{Locating $P_{1\mathrm{dB}}$}{Locating P1dB}}

A coarse sweep brackets the crossing; the crossing is then located precisely.

\begin{implnote}[interpolation versus refinement]
Two estimates are emitted from \emph{one} sweep, so the difference between
methods is a single-variable comparison:
\begin{itemize}
\item \code{p1db\_interpolated\_dbm}: log-linear interpolation between the two
      coarse points bracketing \SI{1}{\dB}.
\item \code{p1db}: bracketed bisection with \emph{additional nonlinear solves}
      inside the bracket (\code{refine\_p1db}), to a tolerance
      \code{--p1db-power-tol-db} (default \SI{0.1}{\dB}).
\end{itemize}
\code{refine\_p1db} validates its bracket first --- it requires
$f(\text{low})<0\le f(\text{high})$ --- and raises rather than returning a
meaningless bisection of a non-bracketing interval.
\end{implnote}

\begin{measured}[refinement versus interpolation]
\begin{center}
\begin{tabular}{@{}l r r r@{}}
\toprule
Device & Published & Refined & $\Delta$ \\
\midrule
JTWPA \SI{6.6}{\giga\hertz} & $-111.458017$ & $-111.118089$ & $+\SI{0.339928}{\dB}$ \\
2c \SI{7.440816}{\giga\hertz} & $-95.083826$ & $-94.857885$ & $+\SI{0.225941}{\dB}$ \\
\bottomrule
\end{tabular}
\end{center}
Each run's interpolated value reproduces the published number to
$\SI{5.1e-9}{\dB}$ (JTWPA) and $\SI{6.0e-9}{\dB}$ (2c), confirming the two
methods are being compared on identical sweeps.

The errors are $0.23$--$\SI{0.34}{\dB}$: above the \SI{0.2}{\dB} scale,
single-signed (refined is always higher, so devices compress \emph{later} than
published), and not grounds for re-running the campaigns wholesale.
\end{measured}

\begin{remark}[an earlier, larger claim was wrong]
A previously reported discrepancy of $+0.461/+\SI{2.920}{\dB}$ is withdrawn. It
was measured on a nine-point coarse grid that no published number used; that
grid placed the 2c first crossing \SI{8.4}{\dB} away from its true value, making
the ``$+2.920$'' a grid artefact rather than a method difference.
\end{remark}

\subsection{Curve metadata}

\begin{implnote}[\code{build\_compression\_curve}]
The curve records more than a single number:
\begin{itemize}
\item \code{first\_1db\_crossing\_dbm} --- the first crossing, found only
      between \emph{adjacent valid} points, so a failed solve cannot create a
      spurious crossing;
\item \code{number\_of\_crossings} --- more than one indicates non-monotone
      compression;
\item \code{nonmonotonic\_compression} --- true if compression ever decreases
      with power;
\item \code{failed\_signal\_power\_dbm} --- the powers at which the solve failed
      or produced non-finite compression.
\end{itemize}
A curve with several crossings, or with failures inside the region of interest,
must not be summarised by its first crossing alone.
\end{implnote}

\subsection{Signal-power continuation}
\label{sec:signal-substep}

Marching up in signal power is itself a continuation problem, with a recovery
ladder analogous to \cref{sec:continuation}:

\begin{enumerate}
\item \textbf{previous} --- warm start from the last converged point;
\item \textbf{secant} --- extrapolate from the last two, in signal current;
\item \textbf{rescaled Floquet} --- rescale the small-signal solution to the new
      amplitude;
\item \textbf{signal substep} --- adaptive continuation along
      \code{signal\_substep} between the last converged current and the target,
      with step sizes set in \si{\dB} of signal power and converted to a
      $\tau$ fraction;
\item \textbf{arclength} --- optional, for turning points.
\end{enumerate}

\begin{implnote}[fallback sizing is deliberately coarse]
The substep's fixed-step fallback is sized like the cold-start path
(\code{\_FALLBACK\_FIXED\_STEPS = 20}), \emph{not} as
$\text{span}/\text{min\_step}$. The latter makes the fallback finest exactly
when adaptive stepping has already failed, scheduling thousands of micro-steps
that only the wall-clock deadline ever stops. Measured before the fix: 6881
steps, $\approx\SI{29}{\hour}$, for one power point.
\end{implnote}

Statuses are distinguished --- \code{VALID\_SOLVED},
\code{SIGNAL\_CONTINUATION\_FAILED}, \code{SIGNAL\_SUBSTEP\_STALL},
\code{DEADLINE}, \code{NONFINITE\_STATE} --- so that a timeout is never reported
as a physical boundary.

\section{Conservation diagnostics}

\subsection{Power balance}
\label{sec:power-balance}

In steady state, power supplied must equal power dissipated. The solver forms
both sides independently.

\emph{Internal:} averaged over the torus,
\begin{equation}
  P_{\mathrm{supplied}} = \left\langle \dot{\xvec}\cdot
    \left(\Cmat\ddot{\xvec} + \Gmat\dot{\xvec} + \Kmat\xvec
          + \Bphi\iJ\right)\right\rangle,
  \qquad
  P_{\mathrm{diss}} = \left\langle\dot{\xvec}\cdot\Gmat\dot{\xvec}\right\rangle.
\end{equation}

\emph{External:} from port power waves. With
\begin{equation}
  a = \frac{V + Z_0 I}{2\sqrt{Z_0}},\qquad
  b = \frac{V - Z_0 I}{2\sqrt{Z_0}},
\end{equation}
the net power into the network at a tone is
$\tfrac12\sum_{\text{ports}}(|a|^{2}-|b|^{2})$.

\begin{implnote}[two boundary subtleties, both once-live bugs]
First, \code{extract\_port\_waves} removes the port resistor's own current
before forming $a,b$: the KCL value is the \emph{Norton source} current, and the
resistor itself carries $V/Z_0$, so the current entering the network is the
difference. Consequently the corresponding shunt loss must \emph{also} be
removed from the dissipation side --- otherwise the two sides describe different
boundaries and the balance can be $O(1)$ wrong.

Second, when a pump-only reference state is supplied, the external quantities
are formed as \emph{differences} from it. A pump-only solution already contains
harmonic conversion (for example $h{=}1 \to h{=}3$), which is energy-conserving
but is not part of signal/idler saturation; leaving it in gives a spurious
residual.
\end{implnote}

\subsection{Manley--Rowe}
\label{sec:manley-rowe}

The Manley--Rowe relations express photon-number bookkeeping: in a lossless
parametric process, the net photon flux summed over the participating modes
vanishes,
\begin{equation}
  \sum_{\mu}\frac{P_\mu}{\omega_\mu} = 0 .
  \label{eq:manley-rowe}
\end{equation}

Two scopes are computed, and \emph{they are not equally valid}.

\paragraph{Conversion scope.} Restricted to $\{(1,0),(1,-1),(1,1)\}$ --- pump,
signal, idler. This is the three-wave conversion channel and
\eqref{eq:manley-rowe} is a genuine invariant for it.

\paragraph{All-tone scope.} Summed over every retained tone.

\begin{caveat}[never gate on the all-tone quantity]
\code{conversion\_manley\_rowe\_rel\_err} is correct, with a real floor of
$\approx\SI{2.2}{\percent}$ from numerical and truncation effects.

\code{all\_tone\_manley\_rowe\_rel\_err} is \textbf{not a valid invariant}. It
peaks at the \emph{smallest} signal power ($0.533$ for JTWPA, $0.500$ for 2c)
--- near-identical across two very different devices, which is the signature of
a fixed factor rather than of physics. The all-tone sum includes pump-harmonic
conversion that is energy-conserving but outside the signal/idler channel, and
it produces a spurious residual of about one half. It is retained as a
diagnostic and must never be used as a pass/fail gate.
\end{caveat}

\begin{implnote}[evaluability floor]
Below a photon scale of $10^{-28}$ the quotient in \eqref{eq:manley-rowe} is
dominated by cancellation and round-off. Such points are marked
\code{manley\_rowe\_evaluable = 0} rather than reported as a large physical
error, and the scale itself (\code{manley\_rowe\_photon\_scale}) is emitted
alongside so the reader can see why.
\end{implnote}

\subsection{Junction diagnostics}

\code{junction\_diagnostics} reports, per branch,
\code{max\_abs\_phase}, \code{max\_abs\_sine} and \code{min\_cosine} over the
torus. The last is the important one: $\min_t\cos(\psi/\rphi)$ going negative
means the junction has passed $\psi/\rphi=\pi/2$ somewhere in the cycle and has
negative differential inductance there --- a saturation mechanism entirely
distinct from pump depletion.

\section{Spatial attribution}
\label{sec:spatial}

\begin{implnote}[\code{spatial\_profiles}]
Branch-resolved amplitudes along the chain: $|\psi_p|$, $|\psi_s|$, $|\psi_i|$
per branch, plus the three-wave phase
\begin{equation}
  \theta_b = 2\arg\psi_{p,b} - \arg\psi_{s,b} - \arg\psi_{i,b},
\end{equation}
unwrapped along the chain, and its gradient
$\Delta k_{\mathrm{eff}} = \mathrm{d}\theta/\mathrm{d}b$ in radians per cell.
This is the local phase mismatch of \eqref{eq:linear-mismatch}, measured rather
than assumed, and it makes visible where along the device the phase matching
degrades.
\end{implnote}

\begin{caveat}[branch index is validated before use as position]
Treating branch index as cell index is only legitimate if $\Bphi$ is an ordered,
connected two-node chain. \code{\_chain\_branch\_nodes} verifies this and raises
otherwise. For one device studied here the nonlinear branches are monotone but
segmented by five intentional node-number gaps: the spatial mapping validates
\emph{order}, not literal node adjacency, and results must be read accordingly.
\end{caveat}

\code{spatial\_profile\_summary} reduces a profile to an overlap integral
\begin{equation}
  \frac{\sum_b |\psi_{p,b}||\psi_{s,b}||\psi_{i,b}|}
       {\|\psi_p\|\,\|\psi_s\|\,\|\psi_i\|}
\end{equation}
and the branch range over which the pump exceeds \SI{10}{\percent} of its peak
--- i.e.\ the effective interaction length.

\section{Stability of the large-signal state}

A converged periodic solution is not necessarily a stable one. Linearising about
the converged large-signal state and examining the Floquet exponents identifies
growing perturbations.

\begin{implnote}[opt-in, and previously wrong]
Stability checking is opt-in (\code{--check-stability}); without it
\code{stability\_status} stays \code{NOT\_CHECKED}. A deep-saturation solution
without the check is not a stability claim.

A previously reported result --- ``stable, dominant exponent
$+\SI{2.87e-6}{\per\second}$'' --- is \textbf{withdrawn}. Two defects:
\begin{enumerate}
\item the stability module passed pump-harmonic keys where
      \code{assemble\_conversion\_matrix} expects a sideband ladder;
\item $\sigma_{\min}$ was owned by a near-DC sideband, giving a value
      bit-identical with the pump on and off up to $\psi/\rphi = 57000$.
\end{enumerate}
The committed tests could not catch either, because they ran on $I_c=0$ --- a
\emph{linear} circuit --- at a zero state. Both are fixed and mutation-verified;
near-DC cases now return \code{INCONCLUSIVE} with a stated reason rather than a
confident \code{STABLE}.
\end{implnote}

\begin{measured}[exponents at the three operating points]
\begin{center}
\begin{tabular}{@{}l r r r@{}}
\toprule
Device & Zero signal & $P_{1\mathrm{dB}}$ & Deepest saturation \\
\midrule
JTWPA & $-\num{3.3785e8}$ (STABLE) & $-\num{1.7638e8}$ (INCONCL.) & $-\num{1.1316e8}$ (INCONCL.) \\
2c    & $-\num{1.1034e-2}$ (STABLE) & $-\num{1.6045e-2}$ (STABLE) & $-\num{1.8755e-2}$ (STABLE) \\
\bottomrule
\end{tabular}
\end{center}
All exponents negative (decaying), and now state-dependent as they must be.
\textbf{No unstable point was found.} The JTWPA \code{INCONCLUSIVE} entries are
a resonance-refinement non-convergence, not the near-DC guard; production points
are clear of that guard (closest sideband \SI{0.52}{\giga\hertz} for JTWPA,
\SI{0.10}{\giga\hertz} for 2c, against thresholds of $\approx\SI{7}{\mega\hertz}$).
\end{measured}

\begin{caveat}[always quote exponents against $\wp$]
A bare exponent in \si{\per\second} is not interpretable. JTWPA's
$|\sigma|/\wp \approx \num{8e-3}$ is real damping; 2c's
$\approx\num{4e-13}$ is numerically marginal. The same numerical value can mean
opposite things on two devices.
\end{caveat}

\section{Results}

\begin{measured}[four-device compression campaign]
\begin{center}
\begin{tabular}{@{}l r r r r@{}}
\toprule
Device & $S$ & Gain [\si{\dB}] & Input $P_{1\mathrm{dB}}$ [\si{\dBm}]
& Depletion at $P_{1\mathrm{dB}}$ [\si{\dB}] \\
\midrule
JPA     & 2  & $13.305144$ & $-132.732343$ & $-3.368085$ \\
JTWPA   & 10 & $27.541036$ & $-111.458017$ & $-0.351776$ \\
FQJTWPA & 6  & $28.534877$ & $-107.386000$ & $-0.598642$ \\
2c      & 10 & $15.916611$ & $\phantom{0}-95.083826$ & $-0.622737$ \\
\bottomrule
\end{tabular}
\end{center}
\end{measured}

\begin{measured}[frequency-resolved]
JTWPA reached $P_{1\mathrm{dB}}$ at 8 of 10 frequencies (minimum
$\SI{-112.280936}{\dBm}$ at \SI{6.888889}{\giga\hertz}); FQJTWPA at 8 of 10
(minimum $\SI{-107.881062}{\dBm}$ at \SI{7.533333}{\giga\hertz}); 2c at 10 of 10
(minimum $\SI{-94.011690}{\dBm}$ at \SI{7.422222}{\giga\hertz}).
\code{NO\_GAIN\_AT\_OPERATING\_POINT} is retained at the dead-band frequencies
rather than reporting a meaningless $P_{1\mathrm{dB}}$.
\end{measured}

\begin{caveat}[every $P_{1\mathrm{dB}}$ here carries an unquantified uncertainty]
The sideband counts in the table above were selected by a procedure that has
since been retired (\cref{chap:validation}), and production-basis
self-convergence has never been measured. Furthermore, the matched basis
constrains $|q|\le1$ at every sideband count; opening $q$ to $\pm2$ cuts
compression by $\SI{44}{\percent}$, which means the published $P_{1\mathrm{dB}}$
values are biased \emph{low}. These numbers should be read as the output of a
specific, documented truncation --- not as converged predictions.
\end{caveat}
